\section{Measurement limits and further experimental design}
\label{app:evaluation}

The configuration and usage measurements in Appendix~\ref{app:methods}
identify the retained client-side record. They do not recover remote
inference hardware, actual peak local resource use, provider invoices,
or all historical sampling settings. The cost estimates are explicitly
rate-based reconstructions; they are not measurements of billing or
energy use. Human oversight outside the logged interaction is likewise
not observable from the trace.

An opportunistic portfolio also limits comparisons. Counting indexed
questions or mentions in a log does not recover a reliable denominator
of substantive attempts. Training-data overlap and prior familiarity
with constructions remain unknown. Without an ablation or a comparison
run, the experiment cannot isolate the effect of persistence, tools,
browsing, model choice or human oversight.

The separate Claude review provides stronger evidence of mathematical
correctness than the initial internal checks. It does not remove these
measurement limitations or determine discovery priority. The following
protocol would make a further study more informative.

\subsection{Designing a comparative experiment}

A follow-up should retain the broad research setting while improving
the measurements that this run could not provide. A practical protocol
would have the following components.

\begin{enumerate}
\item \textbf{Define the unit of evaluation in advance.}
Use a dated set of problems with explicit subparts and a documented
editorial status. Record every substantive attempt, including abandoned
ones, with the reason for stopping. A randomized target subset or
matched target sets would support comparisons that an opportunistic
portfolio cannot.

\item \textbf{Separate discovery from assessment.}
Maintain different fields for statement coverage, proof validity,
computational verification, prior-art status, and outside review.
Freeze proposed answers before specialist assessment, and preserve
later corrections as new versions. Reviewers should receive the
original question and the mathematical argument before seeing aggregate
success claims.

\item \textbf{Record the experimental configuration.}
Save model identifiers, reasoning settings, harness and package
versions, continuation instructions, tool permissions, and retrieval
dates. Capture raw event logs with timestamps and request identifiers,
along with local CPU time and peak memory. Retain response-level usage, including compaction, and separate the
research interval from subsequent writing. Keep usage estimates and
billing records separate until reconciled.

\item \textbf{Treat verification as a research task.}
For a search, specify the exact predicate and its relation to the
question. Include a small direct implementation that checks examples
without reusing the search representation. For a large certificate,
make coverage and the checking rules explicit, retain malformed-input
and negative controls, and distinguish successful generation from
successful verification.

\item \textbf{Use formalization selectively.}
Small algebraic identities, finite certificates, and critical reduction
lemmas are plausible first targets. Formalizing the entire archive at
once would create a large translation project. Any formal account must
also justify that its statement matches the Notebook, including
quantifiers and conventions; a kernel-checked proof of a different
statement would leave the original task unanswered.

\item \textbf{Measure the review burden.}
Record expert time, corrections, changes in status, and which claims
depend on the same imported theorem. Report a short, fully reviewed
core alongside the larger exploratory archive. This would make it
possible to compare output volume with the work needed to establish
reliable mathematical contributions.
\end{enumerate}

